\section{Coefficient spaces and finer fixed-partition estimates}
\label{l:comparisons-section}

The central-isotypic estimate already controls the whole Fourier matrix. This section compares three features of alternative fixed-coset arguments: evaluation in a coefficient space, pointwise character control, and the number of restriction types allowed by the ambient diagram. The first two give weaker bounds but require different intermediate estimates; the branching refinement retains information that the uniform constant $C_u$ suppresses. Throughout this section $X=M_1\sqcup\cdots\sqcup M_b$, $H_i=\Sym(M_i)$ fixes the complement, and $f\ge0$ has mass at most one and satisfies $f(gH_i)\le B/[G:H_i]$. The constants $C_u$ have the meaning fixed in the introduction.

\subsection{Comparison of the information and the conclusions}

The first three rows of Table~\ref{l:comparison-table} compare methods under the same fixed-coset hypothesis; the remaining rows locate transfers with different orientations, quantifiers, or conditional structure. A fixed family of coset bounds controls one common measure. An estimate averaged over omitted sets first needs a choice of partition or a permutation-dependent modification. An entropy inequality valid for every change of law can be applied to a rare coefficient event; an entropy estimate for the reference law alone cannot justify that step. For kernels, small quotient error must be supplemented by control of the full sign multiplicity operators.

\begin{table}[ht]
\centering\small
\renewcommand{\arraystretch}{1.12}
\begin{tabular}{>{\raggedright\arraybackslash}p{.23\textwidth}>{\raggedright\arraybackslash}p{.22\textwidth}>{\raggedright\arraybackslash}p{.22\textwidth}>{\raggedright\arraybackslash}p{.22\textwidth}}
\toprule
Input or method & Orientation and quantifier & Mass treatment & Fourier conclusion\\
\midrule
Central isotypic projection & $gH_i$, every coset of one law & Subprobability allowed & HS squared $bBC_uD^{-1+(u+2)/b}$\\
Coefficient-space projection & Same fixed cosets & Same retained law & HS squared $bB^2C_uD^{-1+(u+2)/b}$\\
Pointwise character bound & Same fixed cosets & Same retained law & HS squared $bB^2C_uD^{-1+(u+4)/b}$\\
Two-sided observations & Image and inverse-image cosets & Two retained factors & Joint-type bound, Section~\ref{l:two-sided-section}\\
Uniform random omitted set & Average marginal TV error & Nearby conditional law & Alternative retained laws; Sections~\ref{l:random-partition-section} and~\ref{l:domains-section}\\
Entropy on a common event & Every law supported on that event & Condition on the event & Tail $CD^{-1/256}$; Theorem~\ref{l:tilted-entropy}\\
Equivariant quotient kernels & Mean rows, or both rows and columns & Subkernels with explicit loss & Product-type HS bounds; Theorems~\ref{l:equivariant-sixteen} and~\ref{l:three-group}\\
Conditional half-products & Inverse-image cap followed by conditional image caps & Product subprobabilities & Operator $\sqrt{22}D^{-1/40}$; Theorem~\ref{l:tabloid-product}\\
\bottomrule
\end{tabular}
\caption{Inputs that lead to different Fourier transfers. The dimension is that of the ambient irreducible except in the kernel product-type estimates. All Hilbert--Schmidt norms are unnormalized.}
\label{l:comparison-table}
\end{table}

\subsection{Evaluation of coefficient functions}

\begin{lemma}[Evaluation in a coefficient space]\label{l:coefficient-evaluation}
Let $H$ be a finite group and let $\Lambda\subseteq\Irr(H)$. If $F$ is a linear combination of matrix coefficients of the types in $\Lambda$, then
\[
 \max_{h\in H}|F(h)|^2\le\left(\sum_{\tau\in\Lambda}(\dim\tau)^2\right)\E_{h\sim U_H}|F(h)|^2.
\]
\end{lemma}
\begin{proof}
Schur orthogonality makes the functions $\sqrt{\dim\tau}\,\tau(h)_{jk}$ an orthonormal basis of this coefficient space. At every $h$, the sum of their squared absolute values is $\sum_{\tau\in\Lambda}(\dim\tau)^2$, because each $\tau(h)$ is unitary. Cauchy--Schwarz for expansion in this basis proves the assertion. Multiple copies of a representation give the same coefficient functions and do not enlarge the space.
\end{proof}

\begin{proposition}[Coefficient-evaluation transfer]\label{l:coefficient-operator}
For every $u>0$ and irreducible $\rho$ of degree $D$,
\[
 \|\widehat f(\rho)\|_{\op}\le\sqrt{bBC_u}\,D^{-1/2+(u+2)/(2b)}.
\]
\end{proposition}
\begin{proof}
Let $P_i$ project onto all $H_i$-types of degree at most $D^{1/b}$. The $P_i$ commute, and $\prod_i(I-P_i)=0$ by the product-degree bound. Thus $Q_i=P_i\prod_{j<i}(I-P_j)$ are pairwise orthogonal projections summing to identity. Write $v_i=Q_iv$.

For fixed $g,w$, the function $h\mapsto\langle w,\rho(gh)v_i\rangle$ on $H_i$ uses only the selected types. Lemma~\ref{l:coefficient-evaluation} gives
\[
 |\langle w,\rho(g)v_i\rangle|^2
 \le C_uD^{(u+2)/b}\E_{h\sim U_{H_i}}|\langle w,\rho(gh)v_i\rangle|^2.
\]
Average over $g$ with mass $f(g)$. Right averaging replaces $f$ by $f*U_{H_i}$, whose density relative to $U_G$ is at most $B$. Schur orthogonality on $G$ gives
\[
 \sum_gf(g)|\langle w,\rho(g)v_i\rangle|^2
 \le BC_uD^{-1+(u+2)/b}\|w\|^2\|v_i\|^2.
\]
Jensen for a measure of mass at most one and $\sum_i\|v_i\|\le\sqrt b\|v\|$ finish the proof. This argument proves the same numerical inequality as Proposition~\ref{l:orbit-span}, using evaluation of functions instead of dimensions of orbit subspaces.
\end{proof}

\subsection{Two more isotypic estimates}

\begin{lemma}[Coefficient and pointwise character estimates]\label{l:isotypic-alternatives}
Let $H\le G$, let $f\ge0$ have mass at most one and coset cap $B/[G:H]$, and let $\tau\in\Irr(H)$ have degree $a$. For every $\rho\in\Irr(G)$,
\begin{align}
 D_\rho\|\widehat f(\rho)P_{\rho,H,\tau}\|_{\HS}^2&\le B^2a^2,\label{l:coefficient-isotypic}\\
 D_\rho\|\widehat f(\rho)P_{\rho,H,\tau}\|_{\HS}^2&\le B^2a^4.\label{l:pointwise-isotypic}
\end{align}
The first bound comes from coefficient projection on each coset; the second comes from a pointwise bound on the character idempotent.
\end{lemma}
\begin{proof}
Put $F=|G|f$. On $xH$, project $h\mapsto F(xh)$ onto the conjugates of the matrix coefficients of $\tau$. Its uniform squared norm is
\[
 a\left\|\E_{h\in H}F(xh)\tau(h)\right\|_{\HS}^2\le B^2a^2,
\]
because the matrix in the norm has operator norm at most the nonnegative function's mean, which is at most $B$. The entries of $\rho(xh)P_{\rho,H,\tau}$ belong to the $\tau$ coefficient space. Replacing $F$ by this projection therefore preserves the tested Fourier matrix. Averaging the bound over cosets and applying full-group Plancherel proves \eqref{l:coefficient-isotypic}.

For the second proof, the character bound $|\chi_\tau(h)|\le a$ gives
\[
 |(f*e_{H,\tau})(g)|\le\frac{a^2}{|H|}f(gH)\le\frac{Ba^2}{|G|}.
\]
The density of this convolution has squared $L^2(U_G)$ norm at most $B^2a^4$. Plancherel gives \eqref{l:pointwise-isotypic}. No bound on the conditional law inside a coset was required in either proof.
\end{proof}

\begin{proposition}[Coefficient and character Hilbert--Schmidt transfers]
\label{l:coefficient-lift}\label{l:coarse-character-lift}
For the disjoint-block data of this section and every $u>0$,
\begin{align}
 \|\widehat f(\rho)\|_{\HS}^2&\le B^2bC_uD^{-1+(u+2)/b},\label{l:coefficient-hs}\\
 \|\widehat f(\rho)\|_{\HS}^2&\le B^2bC_uD^{-1+(u+4)/b}.\label{l:coarse-hs}
\end{align}
\end{proposition}
\begin{proof}
Trace the covering inequality \eqref{l:projection-cover} against $\widehat f^*\widehat f$. Sum the corresponding estimate of Lemma~\ref{l:isotypic-alternatives} over all selected subgroup types. For $j=2,4$,
\[
 \sum_{\dim\tau\le D^{1/b}}(\dim\tau)^j
 \le C_uD^{(u+j)/b}.
\]
This proves both bounds, including every multiplicity.
\end{proof}

Although \eqref{l:central-global} gives a stronger bound, \eqref{l:coarse-hs} records the method that uses only pointwise character control. For $u=2$, large block size makes the complete subgroup reciprocal sum at most three, so cap $B=4$ gives the concrete bound $48bD^{-1+6/b}$.

\begin{proposition}[Low-type central-projector estimates]\label{l:central-projector-operator}
Let $\eta$ be a probability satisfying the coset caps of this section. Let $P_i$ select its $H_i$-types of degree at most $D^{1/b}$ and put $R_i=\sum_{\dim\tau\le D^{1/b}}(\dim\tau)^2$. Then
\[
 \|\widehat\eta(\rho)P_i\|_{\op}\le BD^{-1/2}R_i,
 \qquad D\|\widehat\eta(\rho)P_i\|_{\HS}^2\le B^2R_i.
\]
Consequently
\begin{align*}
 \|\widehat\eta(\rho)\|_{\op}&\le BC_u\sqrt b\,D^{-1/2+(u+2)/b},\\
 \|\widehat\eta(\rho)\|_{\op}&\le B\sqrt{bC_u}\,D^{-1/2+(u+2)/(2b)}.
\end{align*}
\end{proposition}
\begin{proof}
On $H_i$ take the central projector density $k_i(h)=\sum_{\dim\tau\le D^{1/b}}(\dim\tau)\overline{\chi_\tau(h)}$, relative to uniform subgroup measure. Its transform is $P_i$. Pointwise $|k_i|\le R_i$, so convolution of the density $|G|\eta$ with $k_i$ is bounded in absolute value by $BR_i$. Plancherel yields the first displayed bound. Alternatively project onto all selected coefficient spaces on each coset. They are orthogonal; their squared mean sums to at most $B^2R_i$ by the first proof of Lemma~\ref{l:isotypic-alternatives}. This gives the second bound.

Use the orthogonal $Q_i=P_i\prod_{j<i}(I-P_j)$, apply either bound to $Q_iv$, and sum their norms. The inequalities $\sum_i\|Q_iv\|\le\sqrt b\|v\|$ and $R_i\le C_uD^{(u+2)/b}$ give the conclusions.
\end{proof}

\subsection{Counting only types allowed by branching}

Let $p(j)$ again denote the partition number, and define $M(k)=\sum_{j=0}^kp(j)$. For $\lambda\vdash n$ put $k_\lambda=n-\max(\lambda_1,\lambda'_1)$.

\begin{proposition}[Branching refinement]\label{l:branching-hs}
For an irreducible $\rho_\lambda$ of degree $D$, the disjoint-block caps give
\[
 \|\widehat f(\rho_\lambda)\|_{\HS}^2\le BbM(k_\lambda)D^{-1+2/b}.
\]
The blocks may have unequal sizes.
\end{proposition}
\begin{proof}
Young branching says that each diagram $\tau$ occurring under restriction to a block symmetric group is contained in $\lambda$; conjugating a block to the standard subgroup does not change the set of occurring types \cite[Theorem~9.2]{James1978}. Orient a longest row or column of $\lambda$ horizontally, transposing the restricted diagrams too if necessary. Each $\tau$ then has at most $k_\lambda$ boxes below its first row. Given those lower boxes and the block size, its first row is determined. Thus at most $M(k_\lambda)$ types occur, regardless of multiplicity. Apply \eqref{l:central-global} to each selected small type and bound its squared degree by $D^{2/b}$ in \eqref{l:projection-cover}.
\end{proof}

Appendix~\ref{l:transfer-supplement} develops two averages of low-type
projections, a four-factor coefficient decomposition, and explicit
parameter substitutions. Those alternative estimates retain their
proofs, while the tools above suffice for the transfers that follow.
